\begin{afterword}
\summaryhead

The post replies to Richard Chappell's comment on ``Zombies! Zombies?''. The author accepts that ``apparently conceivable'' should be kept apart from ``logically possible'' and declines the label ``Type-A materialism'' as a bundle. To Chappell's complaint that no logical contradiction had been shown, the post gives seven steps: if inward awareness causes what we say about it, as seems empirically likely, then a world with all the causes of our speech contains consciousness, and the zombie world is logically impossible. Like H2O and Twin Earth's XYZ, the word has an empirical referent.

The law that would align an epiphenomenal consciousness with physical talk is called an extra, improbable postulate. Against the claim that a zombie's words are meaningless, the post argues that maps and accuracy can be understood in purely physical terms. It then adds a premise the previous day's argument needed, that a reflective mind checks each of its parts for reliability, and concedes that its etymology of N'Shama was wrong.

\responsehead

The post is a reply to a commenter, and it is candid about itself. It concedes a word (``miraculous''), a missing premise and an etymology, and its seven steps end in a probability, not the logical contradiction Chappell asked for. Its weaknesses are in how it meets the view it answers.

First, the terminology. The post asks for separate words for ``I don't see a contradiction yet'' and ``logically possible.'' Chalmers already had them, prima facie and ideal conceivability, and the zombie argument as Chalmers states it uses the stricter one. The gap the post describes is real, but it is not hidden in the argument under attack.

Second, the position. The post declines Chappell's type-A label, and then rests its answer on an empirical fact about what ``consciousness'' refers to, on the model of water and H2O. In Chalmers's classification that is type-B materialism, and Chalmers's two-dimensional argument was built against this analogy. The post neither names the position nor answers that argument. Chappell's reply two days later, that logical possibility cannot be an empirical question, points at the same gap: the only account the post gives of how an empirical fact yields a logical impossibility is the Twin Earth analogy.

Third, the premises in dispute. Step 4, that ``consciousness'' refers to what makes one say one is aware, is the premise the epiphenomenalist denies, and Chappell called a premise of this kind question-begging. The reply on meaning assumes that accuracy can be judged from the physical system alone. That is the causal theory of knowledge which Chappell's last point, not quoted, said was being assumed. The Texas example shows a check against a territory the checker can observe; an epiphenomenal property gives it nothing to observe. The repaired argument in part (B) measures reliability across a set of worlds that includes the zombie world, without saying why that is the right set.

Fourth, the target. The post treats the zombie argument as an argument for epiphenomenalism. Chappell's comment said that Chalmers does not commit to that view, and Chalmers said so in a comment three days later. Yudkowsky replied that the two theses entail each other. The dispute is real, but the post leaves out Chappell's postscript, and its reader does not learn of it.

On one point the post has support from its opponent. Chalmers calls the fit between experience and reports, on epiphenomenalism, ``something of a lucky coincidence'' that leaves ``a significant burden of proof.'' The post does not take up Chappell's comparison with a universe fit for life.

In short: a candid reply that concedes its slips, but meets Chappell's objections mostly by restating its own premises, and does not engage the replies Chalmers had already given to its central moves.
\end{afterword}
